% =====================================================
% Chapter 9: Troubleshooting and Design Iterations
% General troubleshooting methodology; the actual log of
% issues encountered by the team is project-specific and
% left as a placeholder.
% =====================================================
\chapter{Troubleshooting and Design Iterations}
\label{ch:troubleshooting}

\section{General Troubleshooting Methodology}

A systematic approach to troubleshooting an analog circuit of this type generally proceeds from the supply inward to the signal path \cite{ref1,ref4}:
\begin{enumerate}
    \item \textbf{Verify power}: confirm supply voltages and current draw at each rail before suspecting the signal path.
    \item \textbf{Verify DC bias}: check the quiescent operating point of every active device against the design calculations of Chapter~\ref{ch:circuitdesign}.
    \item \textbf{Isolate stages}: inject a known test signal progressively later in the signal chain (or progressively earlier) to localise which stage is not behaving as expected.
    \item \textbf{Check for oscillation or instability}, particularly around the AGC feedback loop, which is more prone to instability than a simple fixed-gain stage.
    \item \textbf{Re-check component values and orientation} (e.g. electrolytic capacitor polarity, diode/transistor pinout) against the schematic.
    \item \textbf{Document every problem, its diagnosis, and its resolution}, so that design iterations are traceable.
\end{enumerate}

\section{Troubleshooting Log}

Table~\ref{tab:troubleshooting} is structured to record the actual issues encountered during simulation and hardware bring-up, their diagnosis, and their resolution. It is currently empty because this report was prepared before (or without access to) the team's build and debug log.

% =====================================================
% Table: Troubleshooting log
% Empty by design -- no build/debug log entries are available.
% =====================================================
\begin{table}[H]
\centering
\caption{Troubleshooting log.}
\label{tab:troubleshooting}
\begin{tabularx}{\textwidth}{@{}p{2.2cm}XXX@{}}
\toprule
\textbf{Date/stage} & \textbf{Observed problem} & \textbf{Evidence and root cause} & \textbf{Action and outcome} \\
\midrule
--- & --- & --- & --- \\
\bottomrule
\end{tabularx}
\end{table}

\noindent\placeholder{To be completed with each real issue encountered during simulation and hardware bring-up}


%% TODO: Populate the troubleshooting log above with real entries as
%% issues are encountered and resolved during simulation and hardware
%% bring-up. Each entry should include the observed symptom, the
%% diagnostic evidence gathered, the identified root cause, and the
%% corrective action taken.
